\documentclass[11pt,a4paper]{article}

% Preamble for the smc2fc framework documentation.
\usepackage[utf8]{inputenc}
\usepackage[T1]{fontenc}
\usepackage{lmodern}
\usepackage{amsmath, amssymb, amsthm}
\usepackage{algorithm}
\usepackage{algpseudocode}
\usepackage{booktabs}
\usepackage{longtable}
\usepackage{graphicx}
\usepackage{hyperref}
\usepackage{url}
\usepackage{geometry}
\usepackage{xcolor}
\usepackage{listings}
\usepackage{bm}
\usepackage{tikz}
\usetikzlibrary{positioning, arrows.meta, shapes.geometric, fit}

\geometry{margin=1in}

\hypersetup{
    colorlinks=true,
    linkcolor=blue,
    filecolor=magenta,
    urlcolor=cyan,
}

\theoremstyle{definition}
\newtheorem{definition}{Definition}[section]
\newtheorem{proposition}{Proposition}[section]
\newtheorem{theorem}{Theorem}[section]
\newtheorem{remark}{Remark}[section]

% Notation matching old doc so cross-reference works.
\newcommand{\thetadyn}{\bm{\theta}_{\rm dyn}}
\newcommand{\thetactrl}{\bm{\theta}_{\rm ctrl}}
\newcommand{\Lhat}{\widehat{L}}
\newcommand{\xhat}{\widehat{x}}
\newcommand{\dd}{\mathrm{d}}
\newcommand{\E}{\mathbb{E}}
\newcommand{\R}{\mathbb{R}}
\newcommand{\Pcal}{\mathbb{P}}
\newcommand{\Ucal}{\mathcal{U}}
\newcommand{\Xcal}{\mathcal{X}}
\newcommand{\Ncal}{\mathcal{N}}
\newcommand{\Phimax}{\Phi_{\rm max}}
\newcommand{\Fmax}{F_{\rm max}}
\newcommand{\sigmoid}{\mathrm{sigmoid}}
\DeclareMathOperator*{\argmin}{argmin}
\DeclareMathOperator*{\argmax}{argmax}

% Code listing style
\definecolor{codegreen}{rgb}{0,0.6,0}
\definecolor{codegray}{rgb}{0.5,0.5,0.5}
\definecolor{codepurple}{rgb}{0.58,0,0.82}
\definecolor{backcolour}{rgb}{0.97,0.97,0.95}
\lstdefinestyle{code}{
    backgroundcolor=\color{backcolour},
    commentstyle=\color{codegreen},
    keywordstyle=\color{magenta},
    numberstyle=\tiny\color{codegray},
    stringstyle=\color{codepurple},
    basicstyle=\ttfamily\footnotesize,
    breakatwhitespace=false,
    breaklines=true,
    captionpos=b,
    keepspaces=true,
    numbers=left,
    numbersep=5pt,
    showspaces=false,
    showstringspaces=false,
    showtabs=false,
    tabsize=2,
    language=Python,
}
\lstset{style=code}


\title{The \texttt{smc2fc} Framework \\
\large A Practical Reference for Sequential Monte Carlo$^2$ \\
Filtering and Stochastic Optimal Control}
\author{smc2fc development team}
\date{\today}

\begin{document}

\maketitle

\begin{abstract}
\noindent
This document is a practical reference for the \texttt{smc2fc} package
--- a JAX implementation of Sequential Monte Carlo$^2$ for Bayesian
filtering of state-space models with intractable likelihoods, paired
with a control-as-inference receding-horizon controller. It is
deliberately model-independent: the underlying SDE model and observation
likelihood are abstractions supplied by the user, and nothing in this
document depends on the specifics of any particular application.

The intended reader is a research engineer who needs to use, extend, or
tune the framework. The treatment is technical where it must be (the
Bayesian filtering recursion, the unbiased likelihood property, the
adaptive tempering schedule) and pedagogic where it helps (the
control-as-inference framing, the receding-horizon principle, the
soft-barrier feasibility argument). Two kinds of empirical findings are
documented: the JAX-native compile-once architecture (which speeds up
inference by an order of magnitude) and the result that NUTS, despite
being the textbook upgrade from HMC, is roughly 100$\times$ slower than
HMC inside this particular nested control flow. Both are written so
future agents do not repeat the dead-end.

The document deliberately omits speculation about algorithmic
alternatives the framework does not currently implement. The one open
direction it does cover --- improving the soft-HMC controller's wall-clock
without sacrificing control quality --- is the only bottleneck the
authors have hit empirically.
\end{abstract}

\tableofcontents
\newpage

\section{Overview} \label{sec:overview}

The \texttt{smc2fc} package implements two coupled algorithms: a
Sequential Monte Carlo$^2$ (SMC$^2$) Bayesian filter for nonlinear,
non-Gaussian state-space models with intractable marginal likelihoods,
and a stochastic-optimal-control receding-horizon controller built on
the same machinery. The two share the same outer-loop tempered SMC and
the same per-particle differentiable likelihood evaluator, applied to
different log-densities: the Bayes posterior in the filter case, the
control-as-inference posterior in the controller case.

This document describes what is implemented today, with enough
mathematical context to read the source files without an external
textbook open in parallel.

\subsection{Scope and non-scope}

\paragraph{In scope.} The framework code under
\path{smc2fc/}: the core SMC$^2$ engine (\path{smc2fc/core/}), the inner
particle filter (\path{smc2fc/filtering/}), the control-spec abstraction
and the tempered-control loop (\path{smc2fc/control/}). All
mathematical content here applies to any user model that conforms to
the framework's two interfaces (Section~\ref{sec:setting}).

\paragraph{Out of scope.} Specific user models. \texttt{smc2fc} is
designed to be coupled to a user-supplied SDE plus observation model;
the FSA family of physiological models in the sibling
\texttt{FSA\_model\_dev} repository is one such model and is documented
separately. Nothing in this document depends on the FSA model's drift,
state space, or parameters.

\subsection{Architecture in one diagram} \label{sec:architecture}

The framework has four layers, depicted in
Figure~\ref{fig:architecture}. From inside out:

\begin{enumerate}
  \item \textbf{Inner particle filter.} For a fixed parameter vector
        $\thetadyn$, runs an SDE-driven bootstrap particle filter over
        the latent state. Returns an unbiased estimator of the marginal
        observation likelihood $p(y_{1:T} \mid \thetadyn)$.
        Implementation: \path{smc2fc/filtering/gk_dpf_v3_lite.py}.
  \item \textbf{Outer SMC$^2$ over parameters.} Runs adaptive tempered
        SMC over the parameter vector $\thetadyn$, calling the inner
        filter to evaluate likelihoods and using HMC moves between
        tempering steps to rejuvenate the particle cloud.
        Implementations: \path{smc2fc/core/tempered_smc.py} (BlackJAX
        wrapper, kept as a reference) and
        \path{smc2fc/core/jax_native_smc.py} (the production
        compile-once backend; see Section~\ref{sec:jax-native}).
  \item \textbf{Controller.} Reuses the outer SMC$^2$ engine, but with
        the data-likelihood replaced by a control-as-inference
        cost-likelihood $\exp(-\beta J(\thetactrl))$ over a finite-
        dimensional schedule parameter $\thetactrl$. Returns a
        posterior cloud over schedules; the posterior mean is the
        applied control. Implementation: \path{smc2fc/control/}.
  \item \textbf{MPC driver.} Orchestrates the receding-horizon loop:
        for each replan time, advance the plant under the previous
        stride's control, refresh the parameter posterior on the new
        observation window, then call the controller for the next
        stride. The driver is user code; the framework provides the
        ingredients.
\end{enumerate}

\begin{figure}[h]
\centering
\begin{tikzpicture}[
  node distance=8mm,
  every node/.style={font=\small},
  layer/.style={rectangle, rounded corners, draw, minimum height=8mm,
                minimum width=46mm, align=center},
  ext/.style={rectangle, draw, dashed, minimum height=8mm,
              minimum width=46mm, align=center, fill=gray!10},
  arrow/.style={-Stealth, thick},
]
\node[ext]   (model)   {User model: SDE + observation + prior};
\node[layer, below=of model] (pf) {Inner particle filter\\\path{smc2fc/filtering/}};
\node[layer, below=of pf]    (smc2) {Outer SMC$^2$ (HMC moves, adaptive tempering)\\\path{smc2fc/core/}};
\node[layer, below=of smc2]  (ctrl) {Controller (control-as-inference)\\\path{smc2fc/control/}};
\node[ext, below=of ctrl]    (mpc)  {User MPC driver: plant $\leftrightarrow$ filter $\leftrightarrow$ controller};

\draw[arrow] (model) -- node[right] {drift, obs, prior} (pf);
\draw[arrow] (pf)    -- node[right] {$\Lhat_N(\thetadyn)$ (unbiased)} (smc2);
\draw[arrow] (smc2)  -- node[right] {posterior cloud over $\thetadyn$} (ctrl);
\draw[arrow] (ctrl)  -- node[right] {posterior cloud over $\thetactrl$} (mpc);
\draw[arrow, dashed] (mpc.east) to[out=0, in=0, looseness=2] node[right] {observations} (smc2.east);
\end{tikzpicture}
\caption{The four layers of the \texttt{smc2fc} framework. Solid blocks
are framework code; dashed blocks are user code. The dashed arrow on
the right closes the receding-horizon loop: each stride's observations
become the next stride's likelihood input.}
\label{fig:architecture}
\end{figure}

\subsection{Reading order}

The document is structured to be readable straight through, but two
shortcuts are useful.

\paragraph{Mathematician's path.} If you want the mathematical anchors
first, read in this order:
Section~\ref{sec:setting} (the HMM setup),
Section~\ref{sec:pf} (Bayesian filtering recursion and the unbiased
likelihood property), Section~\ref{sec:smc2} (the tempered posterior),
Section~\ref{sec:control} (control-as-inference),
Section~\ref{sec:mpc} (receding-horizon principle).

\paragraph{Engineer's path.} If you want to use or modify the code,
read in this order:
Section~\ref{sec:setting} (the two interfaces you must conform to),
Section~\ref{sec:jax-native} (the production compile-once backend and
how to switch between it and the BlackJAX wrapper),
Section~\ref{sec:config} (every knob in \texttt{SMCConfig}, what each
one does and how to tune it), Section~\ref{sec:control} (the two
chance-constrained cost variants and which one to use),
Appendix~\ref{app:code-api} (one-line API lookup).

\subsection{Conventions}

\begin{itemize}
  \item $\thetadyn$ denotes the dynamics-model parameter vector
        (inferred from data); $\thetactrl$ denotes the schedule
        parameter vector (chosen by the controller). Both live in
        $\R^d$ for some user-defined $d$.
  \item $X_n$ is the latent state at step $n$, $Y_n$ the observation,
        $\Phi_n$ (or generic $U_n$) the exogenous control input at the
        same step.
  \item $N$ is the number of \emph{state} particles (inner PF);
        $M$ is the number of \emph{parameter} particles
        (outer SMC$^2$). Code uses \texttt{n\_pf\_particles} and
        \texttt{n\_smc\_particles} respectively.
  \item File:line citations are checked against the current source at
        the time of writing. If the line numbers drift (e.g.\ after a
        refactor), the function name remains the canonical reference.
\end{itemize}

\subsection{What this document does not promise}

The framework is research code, not a textbook implementation. It has
been tuned against a small number of model classes (the FSA
physiological models, the SWAT hydrological model). Hyperparameter
defaults are reasonable starting points for similar classes; for
genuinely different problems they will need re-tuning, and
Section~\ref{sec:config} is the diagnostic guide for doing so. Where
the framework has hit empirical limits, those limits are documented in
Sections~\ref{sec:jax-native}--\ref{sec:control} and
Section~\ref{sec:future-work}, not glossed.

\section{The generic state-space setting} \label{sec:setting}

This section pins down the abstractions \texttt{smc2fc} expects from the
user's model. Everything in subsequent sections depends only on these
two interfaces.

\subsection{The hidden Markov model in continuous time}

The framework assumes a continuous-time stochastic system observed at
discrete times. The latent state $X(t) \in \R^{n_x}$ evolves under a
controlled It\^o stochastic differential equation
\begin{equation}
\dd X(t) \;=\; b(X(t),\, U(t),\, \thetadyn)\, \dd t
            \;+\; \sigma(X(t),\, \thetadyn)\, \dd W(t),
\label{eq:sde}
\end{equation}
with drift $b: \R^{n_x} \times \R^{n_u} \times \R^d \to \R^{n_x}$,
diffusion $\sigma: \R^{n_x} \times \R^d \to \R^{n_x \times n_x}$
(typically diagonal in framework usage), Brownian motion $W(t)$, an
exogenous control input $U(t) \in \R^{n_u}$, and a static parameter
vector $\thetadyn \in \R^d$. Time is discretised on a uniform grid
$t_n = n \Delta t$, $n = 0, 1, 2, \ldots$, and equation~\eqref{eq:sde}
is integrated step-by-step by Euler--Maruyama:
\begin{equation}
X_{n+1} \;=\; X_n \;+\; b(X_n, U_n, \thetadyn) \Delta t
              \;+\; \sigma(X_n, \thetadyn) \sqrt{\Delta t}\, \xi_{n+1},
              \qquad \xi_{n+1} \sim \Ncal(0, I).
\label{eq:em}
\end{equation}
At each time step $n \ge 1$ the user obtains an observation
$Y_n \in \R^{n_y}$ via the user-supplied likelihood
\begin{equation}
Y_n \mid X_n \;\sim\; p\bigl(\,\cdot \mid X_n,\, C_n,\, \thetadyn\bigr),
\label{eq:obs}
\end{equation}
where $C_n$ is any per-step exogenous regressor (e.g.\ a circadian
forcing, a calendar feature) that the user wants to expose to the
observation model. The framework supports Gaussian, Bernoulli, and
mixed channels in $p(Y_n \mid X_n, \cdot)$ — all log-likelihoods are
summed across channels and any per-bin presence mask supplied by the
user.

The framework does not place semantic constraints on $X_n$, $U_n$, or
$Y_n$. It only requires that the drift, diffusion, observation
likelihood, and prior $p(\thetadyn)$ are all evaluable and
JAX-traceable.

\subsection{The two abstractions the framework consumes} \label{sec:interfaces}

\paragraph{(A) The forward simulator: \texttt{SDEModel}.}
A frozen dataclass that bundles the drift, diffusion, observation
samplers, and per-state physical bounds. Used by the user's
\texttt{StepwisePlant} (Section~\ref{sec:mpc}) for closed-loop
ground-truth simulation, and by the test infrastructure for synthetic
data generation. See \path{smc2fc/simulator/sde_model.py}.

\paragraph{(B) The inference target: \texttt{EstimationModel}.}
A frozen dataclass that bundles the propagate function, the diffusion
function, the observation log-weight function, and configuration of
the parameter prior. The framework's outer SMC$^2$ engine reads only
from this. See \path{smc2fc/estimation_model.py}.

\bigskip
The user supplies one of each. The two are typically thin wrappers
around the same underlying drift and observation code; the split exists
because the simulator side wants numpy-friendly per-step samplers
suitable for a Python loop, while the inference side wants pure-JAX
callables suitable for batching over many particles inside
\texttt{jax.vmap}.

\subsection{What the framework does \emph{not} require}

\begin{itemize}
  \item \emph{No closed-form transition density.} The framework only
        needs to be able to \emph{sample} from
        $p(X_n \mid X_{n-1}, U_n, \thetadyn)$, never to evaluate it
        pointwise. This is what makes it applicable to nonlinear SDEs.
  \item \emph{No tractable marginal likelihood.}
        $p(Y_{1:T} \mid \thetadyn)$ never has to be evaluated exactly;
        the bootstrap particle filter (Section~\ref{sec:pf}) provides
        an unbiased estimator that the outer SMC$^2$ then operates on.
  \item \emph{No Gaussianity.} Neither the transition kernel nor the
        observation likelihood needs to be Gaussian. The HMC moves in
        the outer SMC$^2$ act on the parameter vector $\thetadyn$, not
        on the latent state.
  \item \emph{No separation principle.} The framework does not assume
        the controller can be designed assuming a perfect state
        estimate. The receding-horizon loop of Section~\ref{sec:mpc}
        explicitly threads an SMC$^2$ posterior through the controller
        without separating estimation from control.
\end{itemize}

\subsection{Notational summary}

\begin{center}\small
\begin{tabular}{lll}
\toprule
Symbol & Meaning & Code \\
\midrule
$X_n \in \R^{n_x}$    & Latent state at step $n$                    & \texttt{state} (6D in FSA, 35D in SWAT) \\
$Y_n \in \R^{n_y}$    & Observation at step $n$                    & per-channel \texttt{obs\_value} arrays \\
$U_n,\, \Phi_n$       & Exogenous control input at step $n$        & \texttt{Phi} (sometimes \texttt{U}) \\
$C_n$                 & Per-step exogenous regressor               & \texttt{C} \\
$\thetadyn \in \R^d$  & Static parameter vector (inferred)         & \texttt{u} or \texttt{theta\_dyn} \\
$\thetactrl$          & Control-schedule parameter vector          & \texttt{theta\_ctrl} \\
$N$                   & Number of inner-PF state particles         & \texttt{n\_pf\_particles} \\
$M$                   & Number of outer-SMC parameter particles    & \texttt{n\_smc\_particles} \\
$\lambda$             & Tempering parameter, $0 \to 1$             & \texttt{lam} or \texttt{tempering\_param} \\
$\Lhat_N(\thetadyn)$  & Unbiased likelihood estimator              & \texttt{log\_density(u, ...)} \\
\bottomrule
\end{tabular}
\end{center}

\section{The bootstrap particle filter}\label{sec:pf}

\subsection{Bayesian filtering recursion}

Let $X_n \in \R^3$ denote the latent state $(B,F,A)$ at the $n$th
fifteen-minute time step, and let $Y_n$ denote the corresponding
multi-channel observation. For a fixed value of the static parameter
$\thetadyn$, the FSA-v2 model defines a hidden Markov model with
transition density
\begin{equation}
  X_n \mid X_{n-1} \sim p(\,\cdot\mid X_{n-1},\,\Phi_n,\,\thetadyn)
\end{equation}
implicitly through the Euler--Maruyama discretisation of
\eqref{eq:fsa-B}--\eqref{eq:fsa-A}, and observation density
\begin{equation}
  Y_n \mid X_n \sim p(\,\cdot\mid X_n,\,C(t_n),\,\thetadyn)
\end{equation}
given by the four-channel model \eqref{eq:obs-hr}--\eqref{eq:obs-steps}.
Bayesian filtering is the recursive computation of the conditional
distribution $\pi_n(\dd x_n) := p(\dd x_n \mid Y_{1:n},\,\thetadyn)$. It
factorises into a one-step prediction and a Bayes update:
\begin{align}
  \pi_{n\mid n-1}(\dd x_n)
    &= \int p(\dd x_n \mid x_{n-1},\,\Phi_n,\,\thetadyn)\,
            \pi_{n-1}(\dd x_{n-1}), \label{eq:filter-predict}\\
  \pi_n(\dd x_n)
    &\propto p(Y_n \mid x_n,\,C(t_n),\,\thetadyn)\;\pi_{n\mid n-1}(\dd x_n).
    \label{eq:filter-update}
\end{align}
Equations \eqref{eq:filter-predict}--\eqref{eq:filter-update} have closed
form only when both densities are linear-Gaussian (the Kalman case).
For FSA-v2 they do not, and a Monte Carlo approximation is required.

\subsection{Sequential importance resampling}

The bootstrap particle filter (Gordon, Salmond and Smith, 1993; Doucet
and Johansen, 2011) approximates $\pi_n$ by a weighted empirical measure
\begin{equation}
  \pi_n^N(\dd x_n) \;=\; \sum_{i=1}^{N} W_n^{(i)}\,\delta_{X_n^{(i)}}(\dd x_n),
  \qquad \sum_{i=1}^{N} W_n^{(i)} = 1,
\end{equation}
on a population of $N$ \emph{state particles} $\{X_n^{(i)}\}$. The
recursion has three substeps: (i) propagate each particle through the
dynamics, (ii) reweight by the observation likelihood, (iii) resample to
prevent weight degeneracy. This is sequential importance resampling
(SIR), summarised in Algorithm~\ref{alg:sir}. Resampling is performed
\emph{systematically}: with $u_0 \sim \Ucal[0,1)$, the resampled indices
are obtained by walking the cumulative weight $C(j) = \sum_{i \le j} W_n^{(i)}$
at the offsets $u_k = (k+u_0)/N$ for $k=0,\dots,N-1$. Systematic
resampling is $O(N)$, has lower Monte Carlo variance than multinomial
resampling, and is deterministic given $u_0$.

\begin{algorithm}[H]
\caption{Bootstrap SIR for the FSA-v2 filter, given $\thetadyn$}
\label{alg:sir}
\begin{algorithmic}[1]
\State Sample $X_0^{(i)} \sim p(\dd x_0 \mid \thetadyn)$ for $i=1,\dots,N$;
       set $W_0^{(i)} = 1/N$.
\For{$n = 1,2,\ldots$}
  \State \textbf{Propagate.} For each $i$, draw
         $X_n^{(i)} \sim p(\,\cdot\mid X_{n-1}^{(i)},\,\Phi_n,\,\thetadyn)$
         by one Euler--Maruyama step of \eqref{eq:fsa-B}--\eqref{eq:fsa-A}.
  \State \textbf{Reweight.} Set
         $\tilde W_n^{(i)} = W_{n-1}^{(i)}\,p(Y_n \mid X_n^{(i)},\,C(t_n),\,\thetadyn)$
         and $W_n^{(i)} = \tilde W_n^{(i)}/\sum_j \tilde W_n^{(j)}$.
  \State \textbf{ESS check.} Compute
         $N_{\mathrm{eff}} = 1/\sum_i (W_n^{(i)})^2$. If
         $N_{\mathrm{eff}} < N_{\mathrm{thr}}$, resample
         $\{X_n^{(i)}\}$ systematically by $\{W_n^{(i)}\}$ and reset
         $W_n^{(i)} \leftarrow 1/N$.
\EndFor
\end{algorithmic}
\end{algorithm}

The implementation in
\href{run:../smc2fc/filtering/gk_dpf_v3_lite.py}{\texttt{smc2fc/filtering/gk\_dpf\_v3\_lite.py}}
adds two refinements that we mention but do not develop here. First, when
parameter particles are coupled to the state particles in the SMC${}^2$
outer layer of Section~\ref{sec:smc2}, a Liu--West shrinkage move is
applied with an effective-sample-size dependent bandwidth, slowing
parameter-cloud collapse. Second, when the effective sample size
collapses despite resampling, an optional low-rank Sinkhorn optimal
transport step blends a transport plan into the systematic resample.
Both are diagnostic / numerical safeguards rather than changes to the
mathematical statement of Algorithm~\ref{alg:sir}.

\subsection{The unbiased marginal likelihood}\label{sec:pf-likelihood}

A key property of the bootstrap filter, which makes the SMC${}^2$
algorithm of the next section possible, is that it produces an unbiased
estimator of the marginal likelihood. With the unnormalised weights
$\tilde W_n^{(i)}$ from line~4 of Algorithm~\ref{alg:sir}, define
\begin{equation}
  \Lhat_N(\thetadyn) \;:=\; \prod_{n=1}^{T}
    \biggl(\frac{1}{N}\sum_{i=1}^{N} \tilde W_n^{(i)}\biggr).
  \label{eq:unbiased-likelihood}
\end{equation}
\begin{proposition}[Unbiasedness; Del Moral, 2004]\label{prop:unbiased}
For any $N \ge 1$ and any $T \ge 1$,
\(
   \E\!\bigl[\,\Lhat_N(\thetadyn)\,\bigr]
   = p(Y_{1:T} \mid \thetadyn).
\)
\end{proposition}
A proof is in Del Moral (2004, Thm.~7.4.2). For the present notes the
significance of Proposition~\ref{prop:unbiased} is operational: it allows
us to plug $\Lhat_N(\thetadyn)$ in place of the (intractable) exact
likelihood $p(Y_{1:T}\mid\thetadyn)$ inside any sampler over $\thetadyn$
and obtain a procedure that targets exactly the right posterior. This is
precisely what SMC${}^2$ does.

\section{Optimal-transport rescue inside the filter} \label{sec:ot-rescue}

The Liu--West shrinkage of Section~\ref{sec:liu-west} keeps the inner
state cloud well-behaved most of the time. Occasionally --- under a
particularly informative observation, or when the parameter cloud
visits a low-likelihood region during HMC exploration --- the inner
filter's effective sample size collapses despite the smoothed
resampling. The framework includes a low-rank Sinkhorn optimal-transport
``rescue'' kernel that fires when this happens, blending a transport
plan into the systematic resample by a sigmoid weight on the ESS
ratio. This section explains what it does and why it has the form it
does.

\subsection{When SIR plus Liu--West is not enough}

After the Liu--West shrinkage step of Section~\ref{sec:liu-west}, the
inner state cloud's effective sample size $N_{\rm eff}$ is monitored
once per bin. If it falls below an OT-rescue threshold (default
$N_{\rm eff} / N < 0.05$, controlled by \texttt{ot\_ess\_frac}), the
cloud is judged degenerate enough that systematic resampling will
produce too many duplicates and the filter's ability to track the
state will be lost. The rescue is a soft, sigmoid-blended OT
resampling rather than a hard branch.

The blending weight is
\begin{equation}
w_{\rm OT} \;=\; w_{\rm OT}^{\max} \;\cdot\;
   \sigmoid\!\left(\frac{\text{ess\_thr} - N_{\rm eff}/N}
                          {\text{ess\_temp}}\right),
\label{eq:ot-blend}
\end{equation}
with $w_{\rm OT}^{\max}$ a small cap (default $0.01$, so even at full
saturation only 1\% of the resample comes from OT),
$\text{ess\_thr} = \texttt{ot\_ess\_frac}$, and
$\text{ess\_temp} = \texttt{ot\_temperature}$. The output of the
resampling step at each bin is then
\begin{equation}
X^{\rm out} \;=\; (1 - w_{\rm OT})\, X^{\rm sys+LW}
                \;+\; w_{\rm OT}\, X^{\rm OT},
\end{equation}
i.e.\ a convex combination of the standard systematic+Liu--West output
and the OT-resampled output. Because both inputs are
$\Pcal$-preserving and the weight is in $[0, w_{\rm OT}^{\max}]$, the
combination is also $\Pcal$-preserving.

The continuous-blend formulation matters for the same reason
ESS-scaled bandwidth matters in Section~\ref{sec:liu-west}: the outer
SMC$^2$ HMC kernel (Section~\ref{sec:smc2}) needs the inner-PF
log-likelihood to be a smooth function of $\thetadyn$. A discontinuous
ESS-threshold branch would create gradient discontinuities at every
particle that crosses the threshold, sabotaging the HMC step-size
adaptation.

\subsection{The OT step itself --- low-rank Sinkhorn} \label{sec:sinkhorn}

The OT-rescue kernel takes the current particle cloud
$\{X^{(i)}\}_{i=1}^{N}$ with weights $\{W^{(i)}\}$ and produces a new
cloud whose empirical distribution approximates the same distribution
as systematic resampling, but via a transport plan rather than a
multinomial draw. The transport plan's smoothness with respect to
$\thetadyn$ is what makes the gradient go through.

\paragraph{Step 1 --- low-rank kernel factorisation.}
Pick $r$ anchor particles by index (\texttt{anchor\_idx}, sampled
uniformly without replacement) and form the Gaussian kernel between
all $N$ particles and these $r$ anchors:
\begin{equation}
K_{NR}^{(i, j)} \;=\;
   \exp\!\left(-\frac{\|X^{(i)} - X^{({\rm anchor}\,j)}\|^2}
                       {2\,\epsilon}\right),
\qquad i = 1, \ldots, N;\, j = 1, \ldots, r,
\end{equation}
column-normalised for stability. This is the Nystr\"om approximation
of the full $N \times N$ Gaussian kernel by $K_{NR} K_{NR}^\top$. Cost:
$O(N r)$ in time and memory. Implementation:
\path{smc2fc/filtering/transport_kernel.py:compute_kernel_factor}.

The low-rank form is essential because the full $N \times N$ kernel,
combined with Sinkhorn's $O(N^2)$ matrix-vector products, would
exhaust GPU memory at the rank-$\sim$400 cloud sizes the framework
uses. Empirical default $r = 5$ (\texttt{ot\_rank}); tested up to
$r = 50$ on tighter problems.

\paragraph{Step 2 --- Sinkhorn iterations in scaling form.}
Solve the entropy-regularised optimal transport problem with target
marginals $a$ (current importance weights) and $b$ (uniform weights):
find scalings $u, v \in \R^N_{>0}$ such that
\begin{equation}
\mathrm{diag}(u)\, K_{\rm approx}\, \mathrm{diag}(v)
\quad\text{has marginals}\quad (a,\, b),
\qquad K_{\rm approx} := K_{NR} K_{NR}^\top.
\end{equation}
Sinkhorn alternates the row-marginal and column-marginal projections
\begin{equation}
u^{(\ell+1)} = a / (K_{\rm approx} v^{(\ell)}),
\qquad
v^{(\ell+1)} = b / (K_{\rm approx}^\top u^{(\ell+1)}),
\end{equation}
each iteration costing $O(N r)$ via $K_{\rm approx} v = K_{NR}
(K_{NR}^\top v)$. Default $\ell_{\max} = 2$ (\texttt{ot\_n\_iter}).
Implementation: \path{smc2fc/filtering/sinkhorn.py:sinkhorn_scalings}.

A subtle but important detail: the iteration is implemented via
\texttt{jax.lax.fori\_loop}, not \texttt{lax.scan}. The reason
(documented at the top of \path{sinkhorn.py}) is that \texttt{lax.scan}
inside a particle-filter scan branch causes XLA to unroll all
inner iterations, leading to compilation-time blow-up. \texttt{fori\_loop}
keeps the inner iteration as a real on-device loop.

\paragraph{Step 3 --- barycentric projection.}
The transport plan
$\Gamma = \mathrm{diag}(u)\, K_{\rm approx}\, \mathrm{diag}(v)$
implicitly defines a coupling between source and target particles. The
new particle positions are computed as the barycentric projection
\begin{equation}
X_i^{\rm OT} \;=\;
   \frac{u_i \, [K_{\rm approx} (v \odot X)]_i}
        {u_i \, [K_{\rm approx}\, v]_i}.
\label{eq:barycentric}
\end{equation}
This is a convex combination of source positions, so it preserves any
$[0, 1]^d$ box constraints automatically (no logit / sigmoid needed
for state-bounded models). Implementation:
\path{smc2fc/filtering/project.py:barycentric_projection}.

\subsection{The full OT-rescue API}

The three steps above are combined in
\path{smc2fc/filtering/resample.py:ot_resample_lr}, which is the only
function the inner PF actually calls. Its signature:
\begin{quote}\small
\verb|ot_resample_lr(particles, log_weights, rng_key, stochastic_indices,|\\
\verb|               epsilon=0.5, n_iter=10, rank=50) -> (N, n_states)|
\end{quote}
The \texttt{stochastic\_indices} argument tells the OT kernel which
state coordinates are subject to noise (and thus need rescuing);
deterministic state components pass through unchanged. The function is
fully differentiable through \texttt{jax.grad} --- this is the key
practical feature, distinguishing it from systematic resampling which
has zero gradient with respect to particle positions.

\subsection{Implementation pointer}

\begin{itemize}
  \item \path{smc2fc/filtering/transport_kernel.py} ($\sim$75 lines):
        Nystr\"om kernel factorisation, factor matvec.
  \item \path{smc2fc/filtering/sinkhorn.py} ($\sim$90 lines): low-rank
        Sinkhorn iterations via \texttt{fori\_loop}.
  \item \path{smc2fc/filtering/project.py} ($\sim$45 lines): barycentric
        projection.
  \item \path{smc2fc/filtering/resample.py} ($\sim$100 lines): the
        public \texttt{ot\_resample\_lr} API.
  \item \path{smc2fc/filtering/gk_dpf_v3_lite.py:50} imports
        \texttt{ot\_resample\_lr}; the sigmoid-blend interpolation
        between systematic+LW and OT outputs is in the inner-PF body.
\end{itemize}

\subsection{Practical takeaway}

The OT rescue is a \emph{safety net}, not the main rejuvenation
mechanism. With the default \texttt{ot\_max\_weight = 0.01}, the OT
contribution to the final particle cloud is tiny in normal operation
and only becomes appreciable when the cloud is genuinely degenerate.
Tuning advice: leave the OT-rescue defaults alone unless you see
inner-PF log-likelihood NaNs or sustained $N_{\rm eff} < 5\%$ across
many bins. If you do, the lever is \texttt{ot\_max\_weight} (raise to
0.05--0.10) and \texttt{ot\_n\_iter} (raise to 5--10).

\section{SMC$^2$ over parameters: tempered SMC + HMC} \label{sec:smc2}

The bootstrap filter of Section~\ref{sec:pf} assumes the static
parameter vector $\thetadyn$ is known. In practice $\thetadyn$ is
itself the principal object of inference, with a prior $p(\thetadyn)$
and posterior
\begin{equation}
p(\thetadyn \mid Y_{1:T})
   \;\propto\; p(\thetadyn)\;p(Y_{1:T} \mid \thetadyn).
\end{equation}
A direct sampler is unavailable because $p(Y_{1:T} \mid \thetadyn)$ has
no closed form. The SMC$^2$ algorithm of Chopin, Jacob and
Papaspiliopoulos (2013) replaces this likelihood by the unbiased
estimator $\Lhat_N(\thetadyn)$ from \eqref{eq:unbiased-likelihood} and
runs adaptive tempered SMC over the parameter cloud.

\subsection{The tempered posterior}

For a schedule $0 = \lambda_0 < \lambda_1 < \cdots < \lambda_K = 1$
define the sequence of tempered targets
\begin{equation}
\pi_\lambda(\thetadyn) \;\propto\;
   p(\thetadyn)^{1-\lambda}\,\Lhat_N(\thetadyn)^{\lambda},
   \qquad \lambda \in [0, 1].
\label{eq:tempered-target}
\end{equation}
At $\lambda = 0$ the target is the prior; at $\lambda = 1$ it is, in
expectation, the posterior. The interpolation in between bridges the
two in a way that keeps the importance weights well-behaved. Sampling
from $\pi_{\lambda_K} = \pi_1$ is the goal; the intermediate
$\pi_{\lambda_k}$ are stepping stones.

\subsection{Adaptive tempering: choosing $\lambda_k - \lambda_{k-1}$}

A population of $M$ \emph{parameter particles}
$\{\thetadyn^{(m)}\}_{m=1, \ldots, M}$ (default $M = 256$,
\texttt{n\_smc\_particles}) is propagated through the schedule. At
each step the unnormalised tempering weight is
\begin{equation}
\omega_k^{(m)} \;\propto\;
   \left(\frac{\Lhat_N(\thetadyn^{(m)})}
              {\Lhat_N(\thetadyn^{(m)})_{\rm prev}}\right)^{\!\lambda_k - \lambda_{k-1}},
\end{equation}
and an effective sample size
$M_{\rm eff} = (\sum_m \omega_k^{(m)})^2 / \sum_m (\omega_k^{(m)})^2$
is monitored. The interval $\lambda_k - \lambda_{k-1}$ is chosen
\emph{adaptively}, by bisection, to drive $M_{\rm eff} / M$ to a
prescribed target (\texttt{target\_ess\_frac}, default $0.5$): too
large and the cloud collapses; too small and tempering wastes work.
The bisection is capped at \texttt{max\_lambda\_inc} (default $0.05$
cold-start, $0.10$ bridge) so a single step cannot leap to $\lambda =
1$ in one go.

The on-device implementation of this bisection is
\path{smc2fc/core/jax_native_smc.py:_solve_delta_for_ess}; it runs
inside a \texttt{lax.scan} of fixed depth (30 bisection steps) so the
full tempering loop is one compiled XLA graph. This is part of the
JAX-native compile-once architecture covered in
Section~\ref{sec:jax-native}.

\subsection{The HMC move step}

After resampling at temperature $\lambda_k$, the parameter cloud is
\emph{rejuvenated} by Hamiltonian Monte Carlo moves targeting
$\pi_{\lambda_k}$. Without these moves the cloud loses diversity at
every resample step, eventually collapsing onto a single value.

Each move is a leapfrog-integrated trajectory of length
\texttt{hmc\_num\_leapfrog} (default $8$) at step size
\texttt{hmc\_step\_size} (default $0.025$), accepted/rejected by the
Metropolis ratio. The framework runs \texttt{num\_mcmc\_steps} such
moves per tempering level (default $5$ for cold start, $3$ for bridge
warm-start). Implementation:
\path{smc2fc/core/jax_native_smc.py:_hmc_step_chain}, which is
\texttt{jax.vmap}'d across particles.

\subsection{Adaptive diagonal mass matrix} \label{sec:mass-matrix}

HMC's per-dimension step size is set jointly by \texttt{hmc\_step\_size}
and the inverse mass matrix $M^{-1}$, which acts as a preconditioner.
A poorly-chosen mass matrix collapses the acceptance rate to zero very
quickly during tempering --- for the kinds of high-dimensional
posteriors the framework targets (35-D SWAT, 37-D FSA-v5), an
identity mass matrix collapses HMC at $\lambda \approx 0.3$.

The framework re-estimates a \emph{diagonal} mass matrix at each
tempering level from the current parameter cloud:
\begin{equation}
M^{-1}_{ii} \;=\; \mathrm{Var}_m\!\bigl([\thetadyn^{(m)}]_i\bigr)
                 \,+\, \epsilon, \qquad i = 1, \ldots, d.
\label{eq:diag-mass}
\end{equation}
Implementation: \path{smc2fc/core/mass_matrix.py:estimate_mass_matrix}
(20 lines, deliberately minimal). The regularisation $\epsilon =
10^{-4}$ floors the variance to prevent numerical issues for
parameters that have collapsed to a delta.

A full (non-diagonal) mass matrix would in principle be a better
preconditioner, but on the parameter scales the framework targets the
empirical covariance is too noisy at $M = 256$ particles to invert
stably. The diagonal version is a deliberate engineering choice; see
the comment block in \path{mass_matrix.py} for the empirical history.

\subsection{The full SMC$^2$ loop}

Algorithm~\ref{alg:smc2} stitches the pieces together.

\begin{algorithm}[H]
\caption{Adaptive-tempered SMC$^2$ over $\thetadyn$}
\label{alg:smc2}
\begin{algorithmic}[1]
\State Sample $\thetadyn^{(m)} \sim p(\thetadyn)$ for $m = 1, \ldots, M$;
       set $\lambda_0 = 0$.
\State For each $m$, run the inner bootstrap filter
       (Algorithm~\ref{alg:sir}) with $\thetadyn^{(m)}$ to obtain
       $\Lhat_N(\thetadyn^{(m)})$.
\While{$\lambda_k < 1$}
  \State \textbf{Adaptive $\lambda$.} Solve for the smallest increment
         $\delta \in (0, 1 - \lambda_k]$ such that the tempering ESS
         ratio with weights
         $\omega^{(m)} \propto \Lhat_N(\thetadyn^{(m)})^{\delta}$ equals
         the target. Cap at \texttt{max\_lambda\_inc}. Set
         $\lambda_{k+1} = \lambda_k + \delta$.
  \State \textbf{Reweight \& resample.} Compute weights
         $\omega^{(m)}$, normalise, systematically resample.
  \State \textbf{Re-estimate mass matrix.} $M^{-1}_{ii} =
         \mathrm{Var}_m([\thetadyn^{(m)}]_i) + \epsilon$ from the
         resampled cloud.
  \State \textbf{HMC moves.} Run \texttt{num\_mcmc\_steps} HMC
         trajectories of length \texttt{hmc\_num\_leapfrog} at step
         size \texttt{hmc\_step\_size}, targeting $\pi_{\lambda_{k+1}}$
         on each $\thetadyn^{(m)}$. Re-evaluate
         $\Lhat_N(\thetadyn^{(m)})$ for any moved particle by
         re-running the inner filter.
  \State $k \leftarrow k + 1$.
\EndWhile
\end{algorithmic}
\end{algorithm}

\subsection{Implementation pointer}

Two backends implement Algorithm~\ref{alg:smc2}:
\begin{itemize}
  \item \path{smc2fc/core/tempered_smc.py} --- BlackJAX-wrapped
        version. Easier to read, slower in production because of a
        per-window JIT recompilation cost.
  \item \path{smc2fc/core/jax_native_smc.py} --- production
        compile-once version. Same algorithm, same outputs, but the
        whole tempering loop is one XLA graph that compiles once per
        process and is reused across all subsequent windows.
\end{itemize}
The two are interchangeable from the caller's perspective. Section
\ref{sec:jax-native} explains why the JAX-native version exists, what
it costs to maintain, and what it returns at the cost of a few hundred
extra lines of careful JAX engineering.

\section{JAX-native SMC: once-per-run compilation} \label{sec:jax-native}

This section covers what is, in the authors' view, the framework's
single most consequential engineering decision: the JAX-native
re-implementation of the tempered-SMC outer loop in
\path{smc2fc/core/jax_native_smc.py}, eliminating the per-window JIT
recompilation cost of the BlackJAX-wrapped baseline and reducing the
end-to-end runtime of a multi-window inference run by roughly an
order of magnitude.

It also documents a failed experiment --- replacing HMC with NUTS
(\path{smc2fc/core/jax_native_smc_nuts.py}) --- so that future
contributors do not repeat it.

\subsection{The compile-cost problem with the BlackJAX-wrapped backend}

The reference implementation \path{smc2fc/core/tempered_smc.py} is a
thin wrapper around BlackJAX's
\texttt{blackjax.smc.tempered.build\_kernel}. BlackJAX bakes the
log-likelihood function into the closure of the returned kernel.
That is fine for a single run; the issue is that in the rolling-window
setting (Section~\ref{sec:mpc}) the log-likelihood closes over per-
window observation data, so a new kernel is built per window. Every
fresh kernel triggers a new \texttt{jax.jit} cache entry, and every
new cache entry costs a full HLO recompile --- on the model classes
the framework targets, in the order of $\sim 15$ seconds per window.

For a $27$-window benchmark that recompilation dominates the total
wall-clock; the actual sampling work is only a small fraction of the
visible cost. This is a structural problem of the BlackJAX API, not
of BlackJAX's HMC kernel itself --- the kernel is fine; the issue is
that it is rebuilt every window.

\subsection{The fix}

The fix has three parts, all inside
\path{smc2fc/core/jax_native_smc.py}:

\paragraph{(a) Module-level HMC kernel.} The HMC kernel itself is
built once at module import:
\begin{lstlisting}
HMC_KERNEL = hmc.build_kernel()  # smc2fc/core/jax_native_smc.py:53
HMC_INIT   = hmc.init
\end{lstlisting}
These are module-level constants, not per-call closures. The kernel's
trace cache key depends only on the shape and dtype of its arguments,
not on their values. Same shape/dtype across windows $\Rightarrow$ one
compile per process.

\paragraph{(b) Pytree-Partial wrapping of per-window callables.} The
per-window log-density and the prior log-density are passed in as
\texttt{jax.tree\_util.Partial}-wrapped callables. Partial is a
JAX-recognised pytree, so its bound-arg \emph{shapes/dtypes} enter the
trace cache key but its bound-arg \emph{values} do not. Across
windows the data values change but the data shapes don't, so the
cache hits (\path{jax_native_smc.py:121--130}).

\paragraph{(c) On-device tempering loop via \texttt{lax.while\_loop}.}
The full adaptive-tempering loop, including the ESS bisection that
chooses each $\delta\lambda$, is implemented inside a single
\texttt{lax.while\_loop} (\path{jax_native_smc.py:252}). The
ESS bisection is itself a fixed-iteration \texttt{lax.scan}
(\path{jax_native_smc.py:_solve_delta_for_ess}, 30 bisection steps).
There is no Python-side iteration, no Python-side branching, no
Python-side ESS solver. The whole thing is one XLA graph.

The combined effect: the first window's call pays the JIT compile
cost ($\sim 1.6$ seconds for the cached HMC kernel); every subsequent
window reuses the compiled artefact and pays only the sampling cost
($\sim 1.1$ seconds per window in the benchmark).

\subsection{Empirical speedup}

On a synthetic benchmark target with $256$ parameter particles and
$20$ tempering steps:
\begin{center}
\begin{tabular}{lr}
\toprule
Stage & HMC backend (jax-native) \\
\midrule
JIT compile + first run        & 1.63\,s \\
Cached run (subsequent windows) & 1.09\,s per window \\
\bottomrule
\end{tabular}
\end{center}

For the same benchmark, the BlackJAX-wrapped backend pays roughly
$15$ seconds of recompilation per window (figures will vary by
hardware and HLO cache state); the JAX-native backend amortises this
to a one-shot $1.6$\,s and is then bounded only by the cost of
sampling. The net speedup at $27$ windows is approximately one order
of magnitude.

\subsection{Failed experiment: NUTS as a drop-in replacement} \label{sec:nuts-failed}

The textbook upgrade from HMC is the No-U-Turn Sampler (NUTS), which
adapts the trajectory length per step based on a U-turn diagnostic. A
JAX-native NUTS variant of \path{jax_native_smc.py} was implemented
in \path{smc2fc/core/jax_native_smc_nuts.py} (only the kernel
differs --- everything else is structurally identical). Both kernels
were benchmarked side-by-side on the same synthetic target inside the
same conda environment.

\begin{center}
\begin{tabular}{lrr}
\toprule
Stage & HMC & NUTS \\
\midrule
JIT compile + first run        & 1.63\,s & 102.03\,s \\
Cached run (per window)         & 1.09\,s & 94.23\,s \\
\bottomrule
\end{tabular}
\end{center}

The NUTS version is approximately \emph{60$\times$ slower at compile
time and $85\times$ slower per cached window}.

The cause is structural, not parametric. NUTS dynamically determines
the trajectory length using an internal \texttt{while\_loop} (the
U-turn check). When this dynamic-control-flow loop is placed inside
a \texttt{jax.vmap} (over the $256$ parameter particles) and an outer
\texttt{lax.while\_loop} (the tempering schedule), JAX is forced to
either unroll the dynamic loop --- which produces an enormous XLA
graph --- or compile a deeply nested control-flow structure that the
GPU executes inefficiently. Either way the cost compounds across the
two outer loops.

The trajectory-length expansion is also unbounded in the worst case.
``$5$ NUTS steps per particle per tempering level'' (matching the HMC
default \texttt{num\_mcmc\_steps = 5}) translates into $5 \times 8 =
40$ gradient evaluations under HMC, but anywhere from $5 \times 1 =
5$ to $5 \times 1024 = 5120$ evaluations under NUTS depending on the
maximum tree depth.

\paragraph{Conclusion.} HMC remains the production rejuvenation
kernel. The NUTS implementation file
\path{smc2fc/core/jax_native_smc_nuts.py} is preserved as a documented
dead-end --- not deleted, so the next contributor who has the same
``why don't we just use NUTS?'' idea can read this section, see the
benchmark, and not repeat the experiment.

\subsection{How to switch backends}

In user code, the backends are interchangeable at the import line:

\begin{lstlisting}
# Production: JAX-native, compile-once, HMC.
from smc2fc.core.jax_native_smc import (
    run_smc_window_native, run_smc_window_bridge_native,
)

# Reference: BlackJAX wrapper, slower, easier to read.
from smc2fc.core.tempered_smc import (
    run_smc_window, run_smc_window_bridge,
)

# Documented dead-end (do not use; see Section ref{sec:nuts-failed}).
# from smc2fc.core.jax_native_smc_nuts import (
#     run_smc_window_native, run_smc_window_bridge_native,
# )
\end{lstlisting}

\subsection{Implementation pointers}

\begin{itemize}
  \item \path{smc2fc/core/jax_native_smc.py} ($\sim 430$ lines) ---
        production HMC backend.
  \item \path{smc2fc/core/jax_native_smc_nuts.py} ($\sim 430$ lines)
        --- the dead-end NUTS variant. Read the file's header
        docstring before considering it.
  \item \path{smc2fc/core/mass_matrix.py} ($20$ lines) --- the diagonal
        mass-matrix re-estimation discussed in Section~\ref{sec:mass-matrix}.
  \item \path{smc2fc/core/tempered_smc.py} ($\sim 370$ lines) ---
        BlackJAX-wrapped reference backend.
\end{itemize}

\subsection{Practical takeaway}

For any new application the framework targets, default to
\texttt{jax\_native\_smc.py}. The BlackJAX wrapper is useful only as
a reference when debugging an unfamiliar log-density (its slower
trace cycle gives you cleaner stack traces). The NUTS variant should
not be used.

\section{Gaussian bridge for warm-starts} \label{sec:bridge}

When data arrive sequentially and the framework is run in a rolling-
window mode (Section~\ref{sec:mpc}), the first window is a cold start:
the parameter cloud is sampled from the prior and tempered all the
way to the posterior. Doing the same on every subsequent window would
be wasteful --- successive windows' posteriors are typically close to
each other, so a \emph{warm-start} from the previous window's
posterior is the right thing to do. This section documents the
warm-start mechanism the framework actually uses; it is short by
design.

\subsection{The Gaussian bridge in practice}

The framework's default warm-start mechanism (\texttt{bridge\_type =
"gaussian"}) does the simplest thing that works:

\begin{enumerate}
  \item Fit a single Gaussian
        $q_0 = \Ncal(m_0, \Sigma_0)$ to the previous window's
        posterior cloud by importance-weighted moment match (with
        Liu--West shrinkage on $\Sigma_0$ to reduce its noise).
  \item Run a single inner-PF evaluation on the new window's data per
        previous-cloud particle to get an importance-weighted moment
        match $q_1 = \Ncal(m_1, \Sigma_1)$ of the new posterior.
  \item Initialise the next tempering chain from the blended Gaussian
        $\Ncal(m_b, \Sigma_b)$ with $m_b = (1 - \rho) m_0 + \rho m_1$,
        $\Sigma_b$ a similar interpolation, and $\rho \in [0, 1]$ a
        bias toward the new evidence (default $\rho = 0.7$ in
        \texttt{sf\_blend}, though the field's value is only used by
        the SF variant; see below).
  \item Run the outer SMC$^2$ tempering loop from this initialisation,
        with a slightly relaxed schedule
        (\texttt{max\_lambda\_inc\_bridge = 0.10}, double the cold-
        start cap) and fewer HMC moves per level
        (\texttt{num\_mcmc\_steps\_bridge = 3}, vs $5$ for cold start).
\end{enumerate}

The net effect is that each warm-start window typically converges in
a fraction of the cold-start tempering steps, with a smaller HMC
budget per step. On a $27$-window benchmark this is the difference
between a manageable run and an impractical one.

\subsection{The Schr\"odinger--F\"ollmer / Bures--Wasserstein variant}

\path{smc2fc/core/sf_bridge.py} also implements an alternative
warm-start mechanism based on the closed-form Schr\"odinger--F\"ollmer
bridge between Gaussians, which traces the Bures--Wasserstein geodesic
between $q_0$ and $q_1$ in the manifold of positive-definite
covariance matrices. Activated by \texttt{bridge\_type =
"schrodinger\_follmer"} (with \texttt{sf\_blend = 0.5} for the BW
midpoint), it has the appealing property that its midpoint covariance
is \emph{larger} than either endpoint's when the endpoints differ ---
useful in principle for shedding accumulated bias from previous
windows.

\paragraph{Empirical assessment.} On the model classes the framework
has been tested against, the SF/BW variant has \emph{not} produced
measurably better posteriors than the simple Gaussian bridge above.
The added geometric complexity --- matrix square roots, optional
entropic regularisation, optional info-aware blending keyed off a
local FIM --- is not paying for itself. The plain Gaussian bridge is
therefore the production default and the recommended setting for new
applications.

The SF/BW machinery is preserved in \path{sf_bridge.py} for two
reasons: (a) it may earn its keep on a future, larger-covariance-
mismatch problem; (b) several knobs (\texttt{sf\_info\_aware},
\texttt{sf\_info\_lambda\_thresh\_quantile},
\texttt{sf\_info\_blend\_temperature}) are exposed in
\texttt{SMCConfig} for users who want to experiment. Section
\ref{sec:config} summarises them in the configuration reference.

\subsection{Implementation pointer}

\begin{itemize}
  \item \path{smc2fc/core/sf_bridge.py} --- the bridge implementation.
        The Gaussian-bridge path uses a moment-match helper at the
        top of the file; the SF/BW path uses
        \texttt{bw\_geodesic}, \texttt{bures\_wasserstein\_map}, and
        the larger \texttt{fit\_sf\_base} pipeline.
  \item \path{smc2fc/core/jax_native_smc.py:run_smc_window_bridge_native}
        --- the rolling-window entry point that consumes the bridge
        output and runs the relaxed tempering schedule. Note that
        this entry point currently only supports the SF/BW variant
        (the Gaussian bridge is in the BlackJAX-wrapped backend
        \texttt{run\_smc\_window\_bridge}, in
        \path{smc2fc/core/tempered_smc.py}).
\end{itemize}

\section{Control as inference} \label{sec:control}

Sections~\ref{sec:pf}--\ref{sec:bridge} treat the dynamics parameters
$\thetadyn$ as the unknown and the schedule $U$ as given. This section
reverses the perspective: $\thetadyn$ is held fixed at its posterior
mean (or sampled from the posterior cloud), and the schedule itself
becomes the unknown. The framework solves the resulting stochastic-
optimal-control problem by re-using the SMC$^2$ engine of
Section~\ref{sec:smc2}, with the data-likelihood replaced by an
exponential of the negative cost.

\subsection{The cost functional}

For a fixed dynamics-parameter vector $\thetadyn$, an initial state
$x_0$, and a candidate schedule $U: [0, T] \to \R^{n_u}$, define the
trajectory cost
\begin{equation}
J(U) \;=\; \E_{X \mid U}\!\left[\,
   \int_0^T \ell\bigl(X(t),\, U(t)\bigr)\, \dd t
   \,+\, V_f\bigl(X(T)\bigr)\,\right],
\label{eq:cost}
\end{equation}
where the expectation is over the law of $X$ under the SDE
\eqref{eq:sde} driven by $U$, $\ell$ is a stage cost, and $V_f$ is a
terminal cost. The framework treats $\ell$ and $V_f$ as user-supplied
callables; common choices include negative running reward, soft state
barriers, and effort-quadratic penalties. Concrete examples are given
in Section~\ref{sec:control-cost-variants}.

\subsection{Schedule decoder}

The schedule $U$ is parameterised in finite dimensions by a vector
$\thetactrl \in \R^{d_c}$. The framework's reference decoder is a
Gaussian-RBF basis: with $\{\varphi_a(t)\}_{a=1}^{d_c}$ a fixed basis
on $[0, T]$, a logit bias $c_U$ chosen so that $\thetactrl = 0$
produces a canonical default $U_{\rm def}$, and an upper bound
$U_{\max}$,
\begin{equation}
U(t;\,\thetactrl) \;=\;
   U_{\max} \cdot \sigmoid\!\Biggl(\,c_U
      + \sum_{a=1}^{d_c} (\thetactrl)_a\,\varphi_a(t)\Biggr).
\label{eq:schedule-decoder}
\end{equation}
The sigmoid envelope keeps $U(t; \thetactrl) \in [0, U_{\max}]$ for
every $\thetactrl$; the search is unconstrained in $\R^{d_c}$. Other
decoders (piecewise constant, neural network, etc.) plug in via the
same \texttt{ControlSpec} interface; the framework only requires the
decoder to be a JAX-traceable mapping
$\R^{d_c} \to \R^{n_u \times T}$.

\subsection{Tempered control posterior}

Following the control-as-inference tradition (Kappen, 2005;
Toussaint, 2009; Levine, 2018), the framework defines a probability
measure on $\thetactrl$ by treating the negative cost as an
unnormalised log-likelihood. With inverse temperature $\beta > 0$ and
a Gaussian prior
$\thetactrl \sim \Ncal(0,\, \sigma_{\rm prior}^2 I_{d_c})$, the
\emph{control posterior} is
\begin{equation}
q_\beta(\thetactrl) \;\propto\;
   \exp\bigl(-\beta\,J(\thetactrl)\bigr)\;
   \Ncal(\thetactrl;\,0,\,\sigma_{\rm prior}^2 I_{d_c}),
\label{eq:control-posterior}
\end{equation}
where $J(\thetactrl) := J(U(\,\cdot\,; \thetactrl))$. Sampling from
$q_\beta$ is performed by the SMC$^2$ machinery of
Section~\ref{sec:smc2}, but with the unbiased filter likelihood
$\Lhat_N(\thetadyn)$ replaced by the cost-likelihood
$\exp(-\beta\,J(\thetactrl))$. The inverse temperature $\beta$ is
calibrated automatically from the cost spread of the prior cloud, and
the tempered schedule $\lambda \in [0, 1]$ smoothly interpolates from
prior to control posterior.

\begin{proposition}[Posterior mean as a regularised cost minimiser]
\label{prop:control-mean}
Let $q_\beta$ be the control posterior \eqref{eq:control-posterior}
and write $\bar\theta := \E_{q_\beta}[\thetactrl]$. Then
\begin{equation}
\bar\theta \;=\;
   \argmin_{\theta \in \R^{d_c}}
   \Bigl\{\E_{q_\beta}\bigl[\|\theta - \thetactrl\|_2^2\bigr]\Bigr\},
\label{eq:posterior-mean-l2}
\end{equation}
i.e.\ the posterior mean is the $L^2$ minimiser of squared distance
under the tempered measure. When $J$ is approximately quadratic near
its minimum and the prior is not too informative, $\bar\theta$ sits
in a neighbourhood of the unregularised minimiser
$\argmin_\theta J(\theta)$ of size $O(\beta^{-1/2})$.
\end{proposition}

The framework uses $\bar\theta$, computed as the equally-weighted mean
of the SMC$^2$ parameter cloud at $\lambda = 1$, as the controller
output. The corresponding schedule is then
$U(\,\cdot\,; \bar\theta)$ via \eqref{eq:schedule-decoder}.

\subsection{Hard-indicator vs soft-sigmoid chance-constrained variants}
\label{sec:control-cost-variants}

Many practical control problems are most naturally specified via a
\emph{chance constraint} on a state event rather than via a smooth
running cost. A typical form is
\begin{equation}
\Pcal\bigl[X_t \in \mathcal{B}_{\rm bad}(U_t)\bigr] \;\le\; \alpha,
\qquad \forall\, t \in [0, T],
\label{eq:chance-constraint}
\end{equation}
where $\mathcal{B}_{\rm bad}$ is some user-defined ``bad set'' (e.g.\
the basin of attraction of a collapsed equilibrium). The framework
supports two ways of encoding this constraint inside the
control-posterior cost.

\paragraph{(a) Hard indicator --- pure-IS particle weighting.}
Replace the running cost $\ell$ by the indicator
$\ell^{\rm hard}(X_t, U_t) = \mathbb{1}[X_t \in \mathcal{B}_{\rm bad}(U_t)]$
and aggregate as a per-bin violation rate. Designed for pure-SMC$^2$
controllers that re-weight particles by their empirical violation
rate at each tempering step. Indicator-based; not differentiable in
$\thetactrl$.

\paragraph{(b) Soft sigmoid surrogate --- HMC compatible.}
Replace the indicator by a sigmoid surrogate
\begin{equation}
\ell^{\rm soft}(X_t, U_t) \;=\;
   \sigmoid\!\left(\,\frac{\beta_{\rm chance}\,\bigl(d_{\rm bad}(X_t, U_t)\bigr)}
                          {\rm scale}\,\right),
\label{eq:soft-surrogate}
\end{equation}
where $d_{\rm bad}(X_t, U_t)$ is a signed distance (positive when
$X_t$ is in the bad set, negative otherwise). The scaling parameter
$\beta_{\rm chance}$ controls the sharpness; as
$\beta_{\rm chance} \to \infty$ the surrogate converges to the hard
indicator. Differentiable in $\thetactrl$ via the trajectory; designed
for HMC inner kernels.

\paragraph{Empirical comparison: soft works, hard does not.}
Both variants have been implemented and tested side-by-side on the
chance-constrained control of a representative model. The soft
variant, run with HMC inside the outer SMC$^2$ and a $\beta$ annealed
from $\sim 50$ at low $\lambda$ to $\sim 10^4$ near $\lambda = 1$,
converges to the correct chance-constrained optimum and produces
schedules that respect the constraint at the prescribed rate. The
hard variant --- which matches the framework's importance-weighted
SMC$^2$ structure more naturally, but provides no gradient information
to drive HMC --- does \emph{not} converge to the same optimum in
practice. The empirical recommendation is therefore:

\begin{remark}[Use the soft variant]
For chance-constrained control problems, use the soft sigmoid
surrogate \eqref{eq:soft-surrogate} with HMC. The hard indicator
formulation is preserved in
\path{models/fsa_high_res/control_v5.py:evaluate_chance_constrained_cost_hard}
(in the FSA model dev sandbox repository) for documentation /
regression purposes only --- it is not the production cost function,
and should not be used to drive a deployed controller.
\end{remark}

The detailed numerical comparison and the per-bin behaviour
of both variants is in the FSA model documentation; that is where the
empirical run lives. The framework-level point for this document is
that both code paths exist, they target the same chance constraint
\eqref{eq:chance-constraint}, but only the soft path delivers a
working controller in practice.

\subsection{Implementation pointer}

\begin{itemize}
  \item \path{smc2fc/control/control_spec.py} --- the
        \texttt{ControlSpec} dataclass: cost callable, schedule
        decoder, horizon length, initial state, reference parameters.
  \item \path{smc2fc/control/tempered_smc_loop.py} --- the controller's
        outer loop. Built on the same SMC$^2$ engine as the inference
        side; the only difference is that the log-density it operates
        on is the cost-as-likelihood
        $-\beta\,J(\thetactrl)$ rather than $\log \Lhat_N(\thetadyn)$.
\end{itemize}

\section{Receding-horizon model-predictive control} \label{sec:mpc}

This section closes the loop. The filter of Sections~\ref{sec:pf}--%
\ref{sec:bridge} produces a posterior over $\thetadyn$ on each
observation window; the controller of Section~\ref{sec:control}
produces a posterior over $\thetactrl$ given the current state estimate
and the dynamics-parameter posterior. The MPC driver orchestrates the
two: at each replan time, advance the plant under the previous stride's
control, refresh the filter posterior on the new observation window,
and call the controller for the next stride.

\subsection{Finite-horizon stochastic optimal control}
\label{sec:mpc-finite}

Consider a discrete-time stochastic system on a state $X_k \in \Xcal$
and a control $U_k \in \Ucal$, with transition kernel
$X_{k+1} \sim p(\,\cdot \mid X_k,\, U_k)$, a stage cost
$\ell: \Xcal \times \Ucal \to \R$, and a terminal cost
$V_f: \Xcal \to \R$. The \emph{open-loop} finite-horizon problem of
length $H$ is
\begin{equation}
\min_{u_{0:H-1} \in \Ucal^H}
   \;\E\!\left[\,\sum_{k=0}^{H-1} \ell(X_k, u_k)
                + V_f(X_H)\,\Big|\,X_0 = x\right],
\label{eq:open-loop-fh}
\end{equation}
with the expectation over the joint law of $X_{1:H}$ given the chosen
sequence $u_{0:H-1}$. Denote its optimal value by $V_H(x)$ and a
minimising sequence by $u_{0:H-1}^*(x)$. The mapping
$x \mapsto u_0^*(x)$ is the open-loop optimal first-action policy.

\subsection{The receding-horizon principle} \label{sec:mpc-receding}

In \emph{receding-horizon} or \emph{model-predictive} control,
\eqref{eq:open-loop-fh} is solved repeatedly on a sliding window. At
each replan time $t_n$, given a current state estimate $\xhat_n$:
\begin{enumerate}
  \item Solve \eqref{eq:open-loop-fh} for $u_{0:H-1}^*(\xhat_n)$;
  \item Apply only the first action (or the first stride of length
        $\Delta < H$) to the plant;
  \item Observe new data, update the state estimate to $\xhat_{n+1}$,
        return to step 1.
\end{enumerate}
The closed-loop policy is a function of $\xhat_n$ alone, so the
controller is implementable. Stability of the closed loop is delicate.
Mayne and Rawlings (2000) show that under suitable assumptions on the
horizon length and the terminal cost / terminal set, the closed-loop
state remains in a positively invariant set and converges to a
neighbourhood of a desired equilibrium. The framework relies on the
\emph{long-horizon} variant of that theory: the planning horizon $H$
is chosen long enough relative to the slowest time constant in the
dynamics that the truncation error from neglecting trajectory beyond
$H$ is small.

\subsection{Stochastic + output-feedback variants} \label{sec:mpc-variants}

Three flavours of stochastic MPC are commonly distinguished.

\begin{enumerate}
  \item[(i)] \emph{Certainty-equivalent} MPC. Replace the expectation
        in \eqref{eq:open-loop-fh} by evaluation at a point estimate
        $\hat X = \E[X \mid \xhat_n]$ and solve the resulting
        deterministic problem. Cheap, but discards uncertainty.
  \item[(ii)] \emph{Expected-cost} (stochastic) MPC. Keep the
        expectation, approximate it by Monte Carlo (sampling
        trajectories under the stochastic dynamics) or by scenario
        trees.
  \item[(iii)] \emph{Output-feedback} MPC. When $X_k$ is not directly
        observed, the controller acts on a state estimate $\xhat_n$
        produced by an external estimator. For nonlinear non-Gaussian
        systems no \emph{separation principle} holds in general; the
        joint design of estimator and controller is part of the
        modelling, not a free composition.
\end{enumerate}

The \texttt{smc2fc} pipeline is in classes (ii) and (iii)
simultaneously. The state is unobserved, so a state estimator is
required: this is the rolling-window SMC$^2$ filter of
Sections~\ref{sec:smc2}--\ref{sec:bridge}. The expectation in the
cost \eqref{eq:cost} is kept, evaluated by Monte Carlo on the SDE
under common random numbers, and minimised in $\thetactrl$-space by
the control SMC$^2$ of Section~\ref{sec:control}. The full closed-
loop policy is the composition
\[
\text{(observations)}
   \xrightarrow{\text{rolling SMC$^2$ filter}} \xhat_n
   \xrightarrow{\text{control SMC$^2$, posterior mean}} \bar\theta
   \xrightarrow{\eqref{eq:schedule-decoder}} U(\,\cdot\,; \bar\theta)
   \xrightarrow{\text{plant}} \text{(new observations)}.
\]
There is no separation principle here. The fact that this composition
works at all is an empirical claim, validated on the model classes
the framework has been tested against.

\subsection{The closed-loop pipeline --- generic algorithm}
\label{sec:mpc-pipeline}

Algorithm~\ref{alg:mpc} is the framework-level statement. Specific
applications instantiate it via the bench-driver pattern of
Section~\ref{sec:mpc-driver}.

\begin{algorithm}[H]
\caption{\texttt{smc2fc} closed-loop MPC (generic)}
\label{alg:mpc}
\begin{algorithmic}[1]
\State Initialise the plant; initialise an empty observation buffer.
\For{$n = 1, \ldots, N_w$}
  \State \textbf{Filter update.} If $n = 1$, run cold-start SMC$^2$
         (Algorithm~\ref{alg:smc2}) on the current observation window
         to produce a posterior over $\thetadyn$; otherwise warm-start
         from the previous window's posterior via the Gaussian bridge
         (Section~\ref{sec:bridge}).
  \State \textbf{State extraction.} For a small set of $\thetadyn$
         particles drawn from the posterior, run the bootstrap filter
         (Algorithm~\ref{alg:sir}) forward to the stride endpoint and
         compute the importance-weighted mean of the resulting state
         vector; call this $\xhat_n$.
  \State \textbf{Control plan.} Run the control SMC$^2$ of
         Section~\ref{sec:control}, conditioned on $\xhat_n$ and the
         posterior-mean dynamics parameters $\bar\theta_{\rm dyn}$,
         to obtain a posterior-mean schedule coefficient
         $\bar\theta_{\rm ctrl}$.
  \State \textbf{Apply.} Decode
         $U(\,\cdot\,; \bar\theta_{\rm ctrl})$ via
         \eqref{eq:schedule-decoder} and apply only its first stride
         (length $\Delta_s$) to the plant; collect the resulting
         observations and append to the buffer.
\EndFor
\end{algorithmic}
\end{algorithm}

The horizon over which the controller plans (line 5) is the full
remaining macrocycle, but only the first $\Delta_s$ bins of the plan
are committed (line 6). This is the receding-horizon truncation of
Section~\ref{sec:mpc-receding} applied to whatever model the
framework is driving.

\subsection{Driver pattern: plant, filter, controller} \label{sec:mpc-driver}

Application code instantiates Algorithm~\ref{alg:mpc} via three
artefacts.

\paragraph{(a) The plant interface, \texttt{StepwisePlant}.}
The ``plant'' is whatever object the controller acts upon to produce
new observations. It is a user-supplied class with a single mutating
method:
\begin{quote}\small
\verb|advance(stride_bins, U_schedule)|
$\;\to\;$ \texttt{(trajectory, obs\_dict)}
\end{quote}
A call to \texttt{advance} integrates the SDE for \texttt{stride\_bins}
time steps under the supplied schedule, samples the per-channel
observations, and updates the plant's internal state. Pedagogically,
\texttt{StepwisePlant} is the digital embodiment of the ``real world''
against which the controller is evaluated; mathematically, a single
call samples one transition of the joint state-observation process
under the chosen control over a stride of length $\Delta_s$. The
framework does not provide a generic \texttt{StepwisePlant}; each
user model implements its own (the FSA-v5 model's
\texttt{StepwisePlant} is one example).

\paragraph{(b) The MPC driver script.}
The receding-horizon orchestrator is a user-side script. Its
execution sequence implements Algorithm~\ref{alg:mpc} step by step:
for each window $n$ it
\begin{enumerate}
  \item runs the rolling-window filter by calling
        \texttt{run\_smc\_window\_native} (cold-start, $n=1$) or
        \texttt{run\_smc\_window\_bridge\_native} (warm-start,
        $n \ge 2$), both from
        \path{smc2fc/core/jax_native_smc.py};
  \item extracts the posterior-mean dynamics parameters and the
        smoothed latent state $\xhat_n$ at the end of the window;
  \item assembles a \texttt{ControlSpec} (cost functional, horizon
        length, current state estimate, schedule decoder);
  \item calls the control SMC$^2$ loop in
        \path{smc2fc/control/tempered_smc_loop.py}, which returns a
        posterior-mean schedule coefficient $\bar\theta_{\rm ctrl}$
        and a corresponding $U$ schedule;
  \item calls
        \verb|plant.advance(stride_bins, U_schedule)| to apply only
        the first stride of that schedule;
  \item appends the resulting observations to the data buffer and
        shifts the window.
\end{enumerate}

\paragraph{(c) The posterior-mean handoff between windows.}
The smoothed state $\xhat_n$ extracted in step 2 above is what
initialises the next window's particle filter. Consecutive windows
are therefore connected \emph{at the state level} by a single point
estimate (the importance-weighted smoothed mean) and
\emph{at the parameter level} by the Gaussian bridge of
Section~\ref{sec:bridge}. This is the practical realisation of the
warm-start that gives the rolling SMC$^2$ its computational
tractability.

\subsection{Recursive feasibility and soft barriers}
\label{sec:mpc-feasibility}

State constraints in the cost (e.g.\ a fatigue ceiling, an admissible
operating envelope) are enforced through \emph{soft-barrier} penalties
of the form $\lambda_b \cdot \max(g(X_t) - g_{\max}, 0)^2$ rather than
through hard constraints. The choice is deliberate. Under stochastic
dynamics with state-dependent diffusion and unbounded driving noise,
an exact hard constraint $g(X_t) \le g_{\max}$ cannot be guaranteed
by any controller without making the admissible-control set empty for
some achievable initial conditions; this is the standard
\emph{recursive feasibility} difficulty in stochastic MPC. A soft
penalty preserves recursive feasibility --- there is always a
feasible schedule --- and produces the constraint approximately, with
the violation rate decreasing as $\lambda_b$ is increased.

The framework's chance-constraint formulation of
Section~\ref{sec:control-cost-variants} is the more principled
alternative when the ``bad set'' is precisely defined and a violation-
rate budget is the natural specification. Soft barriers and chance
constraints are not mutually exclusive; both can be combined in the
same cost functional.

\subsection{The flat-cost-surface argument} \label{sec:mpc-robustness}

The controller commits the posterior mean $\bar\theta_{\rm ctrl}$
rather than a sample from $q_\beta$ or its full distribution. One
might worry that this discards useful information, especially when
the filter posterior over $\thetadyn$ is uncertain or biased.
Empirically the policy is nevertheless robust to filter parameter
drift, and the explanation is geometric: near the optimum, the cost
$J(\thetactrl)$ is locally flat in the directions explored by the
SMC$^2$ cloud. A first-order Taylor expansion at the posterior mean
gives, for a perturbed estimate $\bar\theta + \delta$,
\begin{equation}
J(\bar\theta + \delta)
   \;\approx\; J(\bar\theta) + \tfrac{1}{2}\delta^\top H \delta,
   \qquad H = \nabla^2 J(\bar\theta),
\end{equation}
so the additional cost from using $\bar\theta$ in place of an unknown
``oracle'' optimum scales with the Hessian. When $H$ is small in the
directions of greatest filter uncertainty, the closed-loop performance
is correspondingly insensitive to filter accuracy. This is the
mechanism by which a posterior-mean controller is reasonably robust
even when the filter is rate-limited (e.g.\ early in the rolling
schedule, before enough data has accumulated to identify the dynamics
parameters tightly).

\section{Configuration and knobs} \label{sec:config}

This section is the configuration reference for \texttt{SMCConfig}
in \path{smc2fc/core/config.py}. For each field: what it does
mathematically, sensible default ranges, and the diagnostic signal
that says ``increase'' or ``decrease this''. The fields are grouped
into four blocks (outer SMC, inner PF, OT rescue, bridge) following
the order in which they appear in the dataclass.

\subsection{Outer SMC$^2$ (the parameter loop)}

\begin{center}\small
\begin{longtable}{lp{3cm}p{8cm}}
\toprule
Field & Default & Meaning + tuning signal \\
\midrule
\endhead
\texttt{n\_smc\_particles}    & $256$  & Parameter cloud size $M$. Variance of posterior moments scales as $1/\sqrt{M}$. Increase if posterior CIs are too wide; decrease for faster runs. \\
\texttt{target\_ess\_frac}     & $0.5$  & ESS target for adaptive tempering. Lower (e.g.\ $0.3$) = faster but riskier. \\
\texttt{num\_mcmc\_steps}     & $5$    & HMC moves per tempering level (cold start). Increase if posterior samples are correlated within a tempering level. \\
\texttt{max\_lambda\_inc}      & $0.05$ & Cap on $\delta\lambda$ per cold-start tempering step. Smaller $\Rightarrow$ more tempering levels but each one safer. \\
\midrule
\texttt{hmc\_step\_size}      & $0.025$ & HMC leapfrog step size. Decrease if HMC acceptance is low; increase if HMC is wasting compute on tiny moves. Diagnostic: target acceptance rate $\approx 0.7$. \\
\texttt{hmc\_num\_leapfrog}   & $8$    & Leapfrog steps per HMC trajectory. Higher = better mixing per move, more compute per move. \\
\bottomrule
\end{longtable}
\end{center}

\subsection{Inner particle filter}

\begin{center}\small
\begin{longtable}{lp{3cm}p{8cm}}
\toprule
Field & Default & Meaning + tuning signal \\
\midrule
\endhead
\texttt{n\_pf\_particles}     & $400$  & State cloud size $N$ inside the bootstrap filter. Variance of $\Lhat_N(\thetadyn)$ scales as $1/\sqrt{N}$. Increase if $\Lhat_N$ is too noisy (HMC acceptance drops because of noisy gradients). \\
\texttt{bandwidth\_scale}      & $1.0$  & Multiplicative scale on the Silverman bandwidth used by the Liu--West shrinkage of Section~\ref{sec:liu-west}. Increase to $1.5$--$2.0$ if you observe state-cloud collapse. \\
\bottomrule
\end{longtable}
\end{center}

\subsection{OT-rescue (Section~\ref{sec:ot-rescue})}

\begin{center}\small
\begin{longtable}{lp{3cm}p{8cm}}
\toprule
Field & Default & Meaning + tuning signal \\
\midrule
\endhead
\texttt{ot\_ess\_frac}        & $0.05$ & ESS threshold for triggering OT-rescue. Below this fraction, the sigmoid-blend OT-resample fires. Leave at default unless you see sustained $N_{\rm eff} < 5\%$ across many bins. \\
\texttt{ot\_temperature}      & $5.0$  & Sharpness of the sigmoid blend in \eqref{eq:ot-blend}. Larger = sharper transition. \\
\texttt{ot\_max\_weight}      & $0.01$ & Cap on the OT contribution to the final particle cloud. Default 1\% is conservative; raise to $0.05$--$0.10$ if OT-rescue is firing but the cloud still degenerates. \\
\texttt{ot\_rank}             & $5$    & Nystr\"om low-rank approximation rank $r$. Higher = better kernel approximation, more compute. Tested up to $50$. \\
\texttt{ot\_n\_iter}          & $2$    & Sinkhorn iterations. Default 2 is enough for the rescue role; raise to $5$--$10$ for higher OT fidelity. \\
\texttt{ot\_epsilon}          & $0.5$  & Gaussian-kernel bandwidth in OT. Smaller = sharper transport plan, more particle-collapse risk. \\
\bottomrule
\end{longtable}
\end{center}

\subsection{Bridge (Section~\ref{sec:bridge})}

\begin{center}\small
\begin{longtable}{lp{3cm}p{8cm}}
\toprule
Field & Default & Meaning + tuning signal \\
\midrule
\endhead
\texttt{num\_mcmc\_steps\_bridge} & $3$ & HMC moves per tempering level on warm-start (vs $5$ for cold start). The bridge initialisation is closer to the new posterior than the prior, so fewer rejuvenation moves are needed per level. \\
\texttt{max\_lambda\_inc\_bridge}  & $0.10$ & Cap on $\delta\lambda$ per warm-start tempering step (double the cold-start cap). \\
\texttt{bridge\_type}             & \texttt{"gaussian"} & One of \texttt{"gaussian"} (single Gaussian + LW shrinkage; the production default), \texttt{"mog"} (Gaussian mixture), \texttt{"schrodinger\_follmer"} (the BW-geodesic variant; see Section~\ref{sec:bridge}'s empirical assessment). \\
\texttt{bridge\_mog\_components}   & $2$    & Number of components when \texttt{bridge\_type == "mog"}. \\
\texttt{sf\_blend}                 & $0.5$  & SF only: $t \in [0, 1]$ along the BW geodesic. $0$ = previous posterior, $1$ = new-posterior moment match. \\
\texttt{sf\_entropy\_reg}          & $0.0$  & SF only: entropic regularisation; $0$ = exact Bures--Wasserstein. \\
\texttt{sf\_q1\_mode}              & \texttt{"is"} & SF only: how to estimate the new-posterior Gaussian. \texttt{"is"} = single-step importance sampling. \texttt{"annealed"} = K-stage tempered SMC with random-walk Metropolis (more robust in high-D, more expensive). \\
\texttt{sf\_annealed\_n\_stages}   & $3$ & Path B (annealed) only: number of intermediate stages. \\
\texttt{sf\_annealed\_n\_mh\_steps} & $2$ & RW-MH moves per Path B stage. \\
\texttt{sf\_use\_q0\_cov}          & \texttt{False} & SF only: decoupled-mode flag --- use the SF mean update but the previous-posterior covariance. Avoids over-inflation when the new-posterior covariance is MCMC-noisy. \\
\texttt{sf\_info\_aware}           & \texttt{False} & SF only: per-eigenvector info-aware blending using the local Fisher information. Holds the previous-posterior mean in weakly-identified directions across windows. Off by default for bit-identical regression. \\
\texttt{sf\_info\_lambda\_thresh\_quantile} & $0.5$ & SF only (when \texttt{sf\_info\_aware = True}): quantile of FIM eigenvalues used as the soft threshold between ``well-identified'' and ``weakly identified''. \\
\texttt{sf\_info\_blend\_temperature} & $1.0$ & SF only: $\tau$ in the sigmoid threshold. Smaller = sharper. \\
\bottomrule
\end{longtable}
\end{center}

\subsection{Other config dataclasses}

\paragraph{\texttt{RollingConfig}.} Window/stride/dt for the rolling-
window framing. Fields: \texttt{window\_days}, \texttt{stride\_days},
\texttt{dt}, \texttt{n\_substeps}, \texttt{max\_windows}.

\paragraph{\texttt{MissingDataConfig}.} An opinionated default model
of missing-data corruption (rest days, random per-channel dropout,
broken-watch gap). Useful for stress-testing the filter against
realistic data quality issues. Fields: \texttt{dropout\_rate},
\texttt{broken\_watch\_days}, \texttt{rest\_days\_per\_week} plus
the per-channel groupings \texttt{active\_channels},
\texttt{passive\_channels}, \texttt{all\_obs\_channels}.

\subsection{Tuning workflow}

In order, when a new application doesn't converge:
\begin{enumerate}
  \item Diagnose first. Plot per-tempering-level ESS, HMC acceptance
        rate, $\Lhat_N$ variance across particles, and
        per-bin inner-PF ESS.
  \item HMC acceptance below $\sim 0.4$? Reduce
        \texttt{hmc\_step\_size}.
  \item HMC acceptance above $\sim 0.9$? Raise
        \texttt{hmc\_step\_size} (you're under-using each move).
  \item Tempering bisection takes too many steps to find the right
        $\delta\lambda$? Lower \texttt{max\_lambda\_inc} (cap each
        step earlier).
  \item Inner-PF ESS frequently dropping below $5\%$? Raise
        \texttt{n\_pf\_particles} first; if not enough, raise
        \texttt{ot\_max\_weight}.
  \item Posterior moments still noisy after step 4? Raise
        \texttt{n\_smc\_particles}.
\end{enumerate}
Do not change more than one knob at a time, and re-collect the
diagnostics before deciding the next change.

\section{Future work: optimising the soft-HMC controller}
\label{sec:future-work}

The chance-constrained variants of Section~\ref{sec:control-cost-variants}
settled the question of \emph{which} cost surrogate to use: the soft
sigmoid surrogate driven by HMC. This section is about the question
that remains: how to make the soft-HMC controller faster without
sacrificing the control quality that justified the choice in the
first place. The hard-IS variant is not in scope here --- it is a
documented dead-end (Section~\ref{sec:control-cost-variants}, the
``empirical comparison'' remark) and is preserved only for regression
testing.

\subsection{The trade-off as observed today} \label{sec:fw-tradeoff}

Running the soft-HMC controller's tempered SMC$^2$ with high-quality
settings --- more leapfrog steps per HMC trajectory, more HMC moves
per tempering level, finer tempering schedule --- produces better
control plans but costs more wall-clock time. This is, currently, the
framework's main practical bottleneck. Planning latency scales
roughly linearly with HMC work per tempering level:
\[
\text{(wall-clock per plan)}
\;\sim\; n_\lambda \cdot n_{\rm mcmc} \cdot n_{\rm leap}
\cdot \text{(grad eval cost)},
\]
where $n_\lambda$ is the number of tempering levels, $n_{\rm mcmc}$ is
the number of HMC moves per level, and $n_{\rm leap}$ is the leapfrog-
step count per move. All three factors trade quality against speed,
and the right operating point depends on the application's planning-
latency budget. The remaining subsections collect the levers and the
open directions for moving the Pareto frontier.

\subsection{The tunable knobs and what each one costs} \label{sec:fw-knobs}

Four configuration fields are first-line levers. Their definitions and
default values are in the configuration reference (Section~\ref{sec:config});
this subsection is about the quality--speed direction of each.

\begin{center}\small
\begin{longtable}{lp{3.5cm}p{3.5cm}p{4cm}}
\toprule
Knob & Quality direction & Speed direction & Diminishing-return knee \\
\midrule
\endhead
\texttt{hmc\_num\_leapfrog} & Higher = better mixing per trajectory; the proposal travels further per move. & Lower = fewer gradient evaluations per move. & Acceptance plateau around the standard $\tau \approx \pi/(2\omega)$ leapfrog tuning; raising past that wastes compute. \\
\texttt{num\_mcmc\_steps}     & Higher = more rejuvenation per tempering level, lower autocorrelation between tempering levels. & Lower = fewer rejuvenation moves per level. & Per-level autocorrelation typically halves with each doubling up to about $5$--$8$, then flattens. \\
\texttt{target\_ess\_frac}    & Higher (closer to $1$) = smaller $\delta\lambda$ per step, gentler tempering, more reliable convergence. & Lower = larger $\delta\lambda$, fewer tempering levels. & Below $\sim 0.3$ the bisection routinely fails to find a feasible step and the run aborts. \\
\texttt{max\_lambda\_inc}     & Smaller = caps each tempering jump tighter, more steps but each one safer. & Larger = bisection allowed to take bigger leaps. & Larger than the empirical first-window mean of $\delta\lambda$ has no effect; smaller than $\sim 0.01$ produces redundant levels. \\
\bottomrule
\end{longtable}
\end{center}

The four knobs interact. Lowering \texttt{target\_ess\_frac} to cut
the tempering-level count typically requires raising
\texttt{num\_mcmc\_steps} to compensate for the larger $\delta\lambda$
per step --- otherwise the cloud cannot keep up with the temperature
change and HMC acceptance collapses. The combinations that survive
this trade are application-specific.

\subsection{Open directions} \label{sec:fw-directions}

The following are candidate avenues for improving the soft-HMC
controller's Pareto frontier. None has been benchmarked; they are
listed in roughly decreasing order of how cheaply they could be
explored.

\paragraph{Profile first.} Before changing any algorithmic knob, run
\texttt{jax.profiler.start\_trace} around a representative call to
\path{smc2fc/control/tempered_smc_loop.py} and view the trace in
TensorBoard. The wall-clock split between (i) HMC leapfrog gradient
evaluations, (ii) the inner-PF likelihood, (iii) the on-device ESS
bisection, and (iv) host--device transfer latency is the precondition
for any informed knob change. The cheapest improvement may simply be
removing a host sync that nobody noticed.

\paragraph{Adaptive HMC step-size schedule along $\lambda$.} The
existing diagonal-mass-matrix re-estimation
(Section~\ref{sec:mass-matrix}) tightens the geometry per tempering
level, but the step size $\epsilon$ stays fixed. A schedule that
starts with a larger $\epsilon$ at low $\lambda$ (where the target is
nearly the prior, broad and forgiving) and shrinks toward a target
acceptance rate of $\approx 0.7$ at high $\lambda$ (where the
posterior is tight and the geometry is most curved) could reduce the
required $n_{\rm leap}$ at the easy end without sacrificing
acceptance at the hard end. BlackJAX has a dual-averaging adapter
that targets a fixed acceptance probability; wiring it into the
JAX-native loop is straightforward.

\paragraph{Per-tempering-level adaptive leapfrog count.} Currently
\texttt{hmc\_num\_leapfrog} is a single integer applied at every
tempering level. The optimal trajectory length depends on the local
geometry and on how far apart consecutive tempering levels are; a
schedule that tracks the per-level $\delta\lambda$ chosen by the
bisection (longer trajectories when $\delta\lambda$ is large) would
let the controller spend its leapfrog budget where it matters.

\paragraph{Reuse the inference posterior as the controller's prior.}
The controller's tempered SMC$^2$ currently uses a fixed Gaussian
prior $\thetactrl \sim \Ncal(0, \sigma^2 I)$ regardless of what the
filter has learned about $\thetadyn$. When the filter posterior is
informative --- in particular, when its uncertainty has collapsed
along the directions the controller cares about --- starting from a
prior that already encodes that information would shorten the
tempering schedule. The technical question is: what is the right
mapping from the filter's $\thetadyn$ posterior to the controller's
$\thetactrl$ prior? A heuristic such as ``set
$\sigma_{\rm prior}^2$ proportional to the filter's posterior
uncertainty along each direction'' is a starting point; a principled
mapping would require some thought about the joint problem.

\paragraph{Per-bridge HMC budget tuning.}
\texttt{num\_mcmc\_steps\_bridge} is currently $3$ (vs $5$ for cold
start). The bridge initialisation
(Section~\ref{sec:bridge}) is by construction close to the new
posterior, so fewer rejuvenation moves should be sufficient. The open
question is whether $3$ is the right number or whether $1$--$2$ would
do, and whether the answer depends on the magnitude of the
window-to-window posterior shift. A simple sweep on a representative
multi-window benchmark would settle it.

\paragraph{Mixed-precision HMC.} The HMC inner loop is dominated by
matrix--vector products in the gradient evaluation; the outer SMC$^2$
ESS arithmetic and tempering bisection are not. Running HMC in
\texttt{float32} and the rest of the pipeline in \texttt{float64} is
a common $\sim 2\times$ speedup on GPU. The risk is that for
ill-conditioned posterior geometries the lower precision degrades
HMC acceptance enough to require more moves per level, eating the
gain. Worth benchmarking.

\paragraph{Smaller controller particle count with re-warm-starts.}
The controller currently uses the same \texttt{n\_smc\_particles} as
the inference side. The control posterior is typically lower-dimensional
($d_c \approx 8$ vs $d_\theta \approx 16$ for typical applications)
and arguably needs fewer particles to characterise its mean
adequately. A reduction to $128$ or $96$ controller particles, paired
with the next window's bridge warm-start (Section~\ref{sec:bridge}
applied to the controller, not just to the filter), would scale the
per-window controller cost down proportionally. This requires
extending the bridge mechanism to the control posterior, which the
current code does not do.

\subsection{What's not on this list}

A few things deliberately not in the future-work list:

\begin{itemize}
  \item Replacing HMC with NUTS. Documented dead-end; see
        Section~\ref{sec:nuts-failed}.
  \item Replacing the soft surrogate with the hard indicator.
        Documented dead-end; see Section~\ref{sec:control-cost-variants}.
  \item Replacing the Gaussian bridge with the SF/BW variant by
        default. Documented as not paying for itself; see
        Section~\ref{sec:bridge}.
\end{itemize}

These are not open questions. Future contributors should consult the
relevant section before re-attempting.


\appendix
\section{Code-API quick reference} \label{app:code-api}

A one-line description of every public function and class in the
\path{smc2fc/} package, organised by sub-package. The intent is the
``which file is X in?'' lookup --- not a substitute for reading the
docstrings.

\subsection{\texorpdfstring{\texttt{smc2fc/core/}}{smc2fc/core/}}

\begin{description}\small
  \item[\texttt{config.SMCConfig}] Master configuration dataclass for
        outer SMC$^2$, inner PF, OT-rescue, and bridge. See
        Section~\ref{sec:config}.
  \item[\texttt{config.RollingConfig}] Window/stride/dt configuration
        for rolling-window framing.
  \item[\texttt{config.MissingDataConfig}] Default model of
        observation corruption (rest days, dropout, broken-watch
        gaps) used in stress tests.
  \item[\texttt{jax\_native\_smc.run\_smc\_window\_native}]
        Cold-start tempered SMC$^2$ for one observation window.
        JAX-native, compile-once. Production entry point for the
        first window.
  \item[\texttt{jax\_native\_smc.run\_smc\_window\_bridge\_native}]
        Warm-start tempered SMC$^2$ for one observation window using
        the SF/BW bridge initialisation. Production entry point for
        windows $n \ge 2$ when \texttt{bridge\_type =
        "schrodinger\_follmer"}.
  \item[\texttt{jax\_native\_smc.run\_tempered\_chain}] The on-device
        tempered loop itself (single-window, single
        \texttt{lax.while\_loop}). Internal but exposed.
  \item[\texttt{jax\_native\_smc.TemperedState}] \texttt{NamedTuple}
        carrying particles, log-weights, current $\lambda$, and HMC
        kernel state across iterations of the on-device loop.
  \item[\texttt{jax\_native\_smc\_nuts.*}] NUTS variant of the above.
        Documented dead-end; see Section~\ref{sec:nuts-failed}. Not
        for production use.
  \item[\texttt{tempered\_smc.run\_smc\_window}] BlackJAX-wrapped
        cold-start equivalent of
        \texttt{run\_smc\_window\_native}. Slower, easier to debug.
  \item[\texttt{tempered\_smc.run\_smc\_window\_bridge}] BlackJAX-
        wrapped Gaussian-bridge warm-start. The
        \texttt{bridge\_type = "gaussian"} path lives here; the
        JAX-native side currently only implements the SF variant.
  \item[\texttt{mass\_matrix.estimate\_mass\_matrix}] Diagonal mass-
        matrix estimator from a particle cloud (variance with a
        small regularisation floor). Re-called per tempering level.
  \item[\texttt{sampling.sample\_from\_prior}] Convenience prior
        sampler in unconstrained space, used at cold start.
  \item[\texttt{sf\_bridge.bures\_wasserstein\_map}] BW transport
        matrix between two PSD covariances.
  \item[\texttt{sf\_bridge.bw\_geodesic}] Closed-form Bures--
        Wasserstein geodesic between two Gaussians, parameterised by
        $t \in [0, 1]$.
  \item[\texttt{sf\_bridge.estimate\_target\_gaussian}] Single-step
        importance-sampling fit of $q_1$ (the new posterior) given
        $q_0$ (the previous posterior).
  \item[\texttt{sf\_bridge.estimate\_target\_gaussian\_annealed}]
        $K$-stage tempered SMC fit of $q_1$ with random-walk
        Metropolis. Activated by \texttt{sf\_q1\_mode = "annealed"}.
  \item[\texttt{sf\_bridge.fit\_sf\_base}] Top-level entry point for
        the SF/BW bridge: returns the blended Gaussian to seed the
        next window.
  \item[\texttt{sf\_bridge.estimate\_fim\_hessian}] Diagonal Fisher-
        information estimate from the new-window log-likelihood
        Hessian; used by the info-aware blending mode.
  \item[\texttt{sf\_bridge.per\_eigenvector\_blend} /
        \texttt{info\_aware\_bridge\_mean}] Per-direction blending
        between $m_0$ and $m_1$ keyed off the FIM eigenvalues.
        Activated by \texttt{sf\_info\_aware = True}.
\end{description}

\subsection{\texorpdfstring{\texttt{smc2fc/filtering/}}{smc2fc/filtering/}}

\begin{description}\small
  \item[\texttt{gk\_dpf\_v3\_lite.make\_gk\_dpf\_v3\_lite\_log\_density\_compileonce}]
        Build the closed-over inner-PF log-likelihood for a single
        observation window. The compile-once variant the JAX-native
        backend consumes.
  \item[\texttt{gk\_dpf\_v3\_lite.make\_gk\_dpf\_v3\_lite\_log\_density}]
        Per-call variant; closes over data values rather than just
        shapes. Used by the BlackJAX-wrapped backend.
  \item[\texttt{\_gk\_kernel.compute\_ess}] Effective sample size from
        log-weights.
  \item[\texttt{\_gk\_kernel.silverman\_bandwidth}] Silverman's rule
        bandwidth for the smoothing kernel.
  \item[\texttt{\_gk\_kernel.log\_kernel\_matrix}] Log of the
        Gaussian kernel matrix between particles, used by the
        Liu--West smoothing step.
  \item[\texttt{\_gk\_kernel.smooth\_resample\_basic}] Plain
        systematic resample with Liu--West shrinkage and Gaussian
        smoothing.
  \item[\texttt{\_gk\_kernel.smooth\_resample\_ess\_scaled}] As
        above, with bandwidth scaled by ESS to avoid over-smoothing
        when ESS is high.
  \item[\texttt{\_gk\_kernel.smooth\_resample}] Adaptive wrapper that
        dispatches to one of the variants above depending on ESS.
  \item[\texttt{\_gk\_kernel.smooth\_resample\_ess\_scaled\_lw}]
        Log-weighted variant for use inside the bootstrap filter
        likelihood evaluation.
  \item[\texttt{transport\_kernel.compute\_kernel\_factor}] Nystr\"om
        low-rank kernel factor between $N$ particles and $r$ random
        anchors; the $K_{NR}$ matrix in
        Section~\ref{sec:ot-rescue}.
  \item[\texttt{transport\_kernel.factor\_matvec}] Mat-vec via the
        Nystr\"om factorisation; $O(Nr)$ instead of $O(N^2)$.
  \item[\texttt{transport\_kernel.factor\_matvec\_batch}] Batched
        version of the same.
  \item[\texttt{sinkhorn.sinkhorn\_scalings}] Sinkhorn iterations on
        the low-rank-kernel-factored transport problem, returning the
        $(u, v)$ scaling vectors.
  \item[\texttt{project.barycentric\_projection}] Barycentric
        projection of a source cloud onto a target cloud given
        Sinkhorn scalings $(u, v)$ and the kernel factor.
  \item[\texttt{resample.ot\_resample\_lr}] Top-level OT-rescue:
        Sinkhorn $\to$ projection $\to$ sigmoid blend with the
        systematic-resample output.
\end{description}

\subsection{\texorpdfstring{\texttt{smc2fc/control/}}{smc2fc/control/}}

\begin{description}\small
  \item[\texttt{control\_spec.ControlSpec}] Dataclass describing the
        controller's task: cost callable, schedule decoder, horizon
        length, initial state estimate, reference dynamics
        parameters.
  \item[\texttt{config.SMCControlConfig}] Configuration dataclass for
        the controller's tempered-SMC$^2$ loop (counterpart to
        \texttt{SMCConfig} on the inference side).
  \item[\texttt{tempered\_smc\_loop.run\_tempered\_smc\_loop\_native}]
        JAX-native control SMC$^2$ loop: tempers the prior toward
        the cost-as-likelihood and returns the posterior cloud
        (mean = recommended $\bar\theta_{\rm ctrl}$).
  \item[\texttt{tempered\_smc\_loop.run\_tempered\_smc\_loop}]
        BlackJAX-wrapped reference equivalent.
  \item[\texttt{rbf\_schedules.RBFSchedule}] Gaussian-RBF schedule
        decoder \eqref{eq:schedule-decoder}: maps a coefficient
        vector to a $U(t)$ trajectory.
  \item[\texttt{calibration.calibrate\_beta\_max}] Choose the high-
        $\lambda$ inverse temperature from the cost spread of the
        prior cloud.
  \item[\texttt{calibration.build\_crn\_noise\_grids}] Pre-sample
        common-random-number diffusion shocks for use across the
        tempering schedule.
  \item[\texttt{diagnostics.plot\_cost\_histogram} /
        \texttt{plot\_schedule\_comparison} /
        \texttt{plot\_trajectories} / \texttt{print\_smc\_step} /
        \texttt{evaluate\_gates}] Post-hoc inspection helpers; not on
        the run-time critical path.
  \item[\texttt{lqg.controller.LQGSpec} / \texttt{LQGController} /
        \texttt{build\_lqg\_open\_loop\_schedule}] LQG baseline with
        a finite-horizon Riccati solver. Provides a sanity-check
        comparator to the tempered-SMC$^2$ controller.
  \item[\texttt{lqg.linearize.linearize\_drift\_at}] Local
        linearisation of the SDE drift at an operating point; used
        by the LQG baseline.
  \item[\texttt{lqg.riccati.solve\_riccati\_backward} /
        \texttt{compute\_lqr\_gain}] Backward Riccati integration and
        LQR gain extraction.
\end{description}

\subsection{\texorpdfstring{\texttt{smc2fc/simulator/}}{smc2fc/simulator/}}

\begin{description}\small
  \item[\texttt{sde\_model.StateSpec}] Description of a single SDE
        state component (name, prior, bounds).
  \item[\texttt{sde\_model.ChannelSpec}] Description of a single
        observation channel (which states it sees, noise model).
  \item[\texttt{sde\_model.SDEModel}] The user-supplied model
        interface: drift, diffusion, observation kernel, parameter
        priors. Section~\ref{sec:setting} expands on this.
  \item[\texttt{sde\_solver\_diffrax.solve\_deterministic\_jax}]
        Deterministic-mean trajectory via diffrax (used for warm-
        starts and diagnostics).
  \item[\texttt{sde\_solver\_diffrax.solve\_sde\_jax}] Stochastic SDE
        solve via diffrax with optional control schedule.
  \item[\texttt{sde\_observations.generate\_all\_channels}] Sample
        per-channel observations from a trajectory.
\end{description}

\subsection{\texorpdfstring{\texttt{smc2fc/transforms/} and top-level}{smc2fc/transforms/ and top-level}}

\begin{description}\small
  \item[\texttt{transforms.unconstrained.build\_transform\_arrays}]
        Build the bound/scale arrays for the unconstrained-space
        reparameterisation.
  \item[\texttt{transforms.unconstrained.constrained\_to\_unconstrained}
        / \texttt{unconstrained\_to\_constrained}] Bijection between
        the user-facing constrained $\thetadyn$ and the HMC-friendly
        unconstrained $u$.
  \item[\texttt{transforms.unconstrained.log\_prior\_unconstrained}]
        Log-prior in unconstrained space (with Jacobian correction).
  \item[\texttt{transforms.unconstrained.split\_theta}] Slice an
        unconstrained vector into dynamics and control halves.
  \item[\texttt{estimation\_model.EstimationModel}] User-side
        inference interface (parameter list, bounds, log-prior,
        link to \texttt{SDEModel}).
\end{description}


\end{document}
